\begin{problem}[The double line]
Let
\[
A=k[x,y]/(y^2).
\]
Determine $\Nil(A)$, $A_{\mathrm{red}}$, the prime ideals of $A$ in terms of primes of
$k[x]$, and the underlying topological support.  Compare with
\[
B=k[x,y]/(y).
\]
Show that $A\not\cong B$ although their spectra are naturally homeomorphic.
\end{problem}
\paragraph{Why this problem matters.}
The double line is the first geometric example in which an entire positive-dimensional
subspace is thickened without adding points.
\paragraph{Useful ideas before starting.}
Every prime contains $\bar y$, and quotienting by $(\bar y)$ gives $k[x]$.
\paragraph{Guided hints.}
Use quotient correspondence.  To prove the rings differ, compare nilpotents.
\begin{solution}
The class $\bar y$ is nilpotent and every nilpotent lies in all primes.  Since
\[
A/(\bar y)\cong k[x]
\]
is reduced, $\Nil(A)=(\bar y)$.  Therefore
\[
A_{\mathrm{red}}\cong k[x].
\]
Prime ideals of $A$ are precisely the inverse images of primes of $k[x]$ under
$A\twoheadrightarrow A/(\bar y)$.  Hence $\Spec A$ is naturally homeomorphic to
$\Spec k[x]$, which is also $\Spec B$.  But $A$ contains the nonzero nilpotent $\bar y$,
whereas $B\cong k[x]$ is reduced, so the rings cannot be isomorphic.
\end{solution}
\paragraph{What happened mathematically.}
The ordinary line and double line share every point and every specialization relation, but
not the same functions.
\paragraph{Extensions.}
Analyze $k[x,y]/(y^n)$ for arbitrary $n$.
\paragraph{Provenance.}
New canonical problem elaborating the legacy doubled-point phenomenon in CA13.
